\documentclass[12pt]{article}

\usepackage[a4paper, margin=1in]{geometry}
\usepackage{longtable}
\usepackage{titlesec}

\title{\noindent\rule{\textwidth}{0.5pt}
\huge{AirWatch} \vspace{0.1cm}\\ \large{Weather \& Air Pollution Monitor\\}
\noindent\rule{\textwidth}{0.5pt}
\Large{Project Snapshot Overview\\\small{All Four Development Snapshots\\}}
\vspace{0.1cm}
\small{ 
Version: 4.0\\ 
Date: May 2026\\
Platform: Mobile (iOS/Android) \& Web\\
CS3338 Group  5}
}
\date{}
\begin{document}

\maketitle

\tableofcontents
\newpage

\section{Project Introduction}
AirWatch is a cross-platform weather and air pollution monitoring application targeting both mobile (iOS/Android) and web browsers. The application aggregates real-time air quality index (AQI) data alongside traditional meteorological readings, providing users with actionable health advisories, historical trend charts, and location-based alerts. The project is simulated across four development snapshots spanning eight months of a typical agile software lifecycle.

\subsection{Snapshot 1 - Project Start}

Snapshot 1 summarizes the initial development phase of the AirWatch project. 
The purpose of this phase is to establish the project foundation, define the 
software architecture, initialize development workflows, and create a working 
prototype capable of displaying weather and air quality information as part of the Senior Design/CS curriculum.

\subsection{Primary Objectives}

\begin{enumerate}
    \item Establish GitHub repository structure
    \item Configure CI/CD pipeline using GitHub Actions
    \item Set up React Native mobile application
    \item Set up React web application
    \item Configure Node.js and Express backend server
    \item Configure PostgreSQL database
    \item Integrate OpenWeatherMap API
    \item Integrate AirVisual API
    \item Build prototype UI screens
    \item Create initial unit tests
\end{enumerate}

\subsection{Frameworks and Dependencies}

\begin{longtable}{|p{5cm}|p{3cm}|p{5cm}|}
\hline
\textbf{Tool} & \textbf{Version} & \textbf{Purpose} \\
\hline
React Native & 0.73 & Mobile application framework \\
\hline
React & 18.2 & Web frontend framework \\
\hline
Node.js + Express & 20 LTS & Backend REST API \\
\hline
PostgreSQL & 15 & Database management \\
\hline
Redux Toolkit & 1.9 & State management \\
\hline
Axios & 1.6 & HTTP requests \\
\hline
Jest & 29 & Testing framework \\
\hline
\end{longtable}

\subsection{Team Roles \&  Task Delegation}

\begin{longtable}{|p{3cm}|p{3cm}|p{7cm}|}
\hline
\textbf{Team Member} & \textbf{Role} & \textbf{Primary Responsibilities} \\
\hline
Brandon Khuu & Project Manager & Sprint planning, risk management, stakeholder updates \\
\hline
Tommy Pang & Lead Mobile \& Web Dev & React Native setup, navigation, location services \newline
React web app, responsive layout, PWA config\\
\hline
Nathan Lopez & Back-End Dev & Express API, database schema, third-party integrations \\
\hline
Joshua Xotoy & QA Engineer & Test plans, manual testing, TestRail management\\
\hline
\end{longtable}



\subsection{Sprint Goals}

\subsubsection{Sprint 1}
\begin{itemize}
    \item Environment setup
    \item Repository initialization
    \item Backend scaffolding
    \item Database connection
\end{itemize}

\subsubsection{Sprint 2}
\begin{itemize}
    \item Weather API integration
    \item AQI API integration
    \item Mobile home screen
    \item Web home page
    \item Initial unit testing
\end{itemize}

\subsection{Workflow}
Sprint Length: 2 weeks  |  Methodology: Scrum  |  Stand-ups: Daily 9 AM
Sprint 1: Setup \& scaffolding → Sprint 2: Core data integration \& prototype screens
Version Control: GitHub Flow – feature branches → PR review → merge to main



\subsection{Expected Deliverables}

\begin{itemize}
    \item Functional prototype
    \item Basic UI implementation
    \item Connected backend API
    \item Database schema
    \item Initial documentation
\end{itemize}
\newpage

\section{Snapshot 2 - Checkpoint 1 Reflection}

Snapshot 1 successfully established the project foundation, backend API,
weather integrations, and prototype user interface.

\subsection{New Feature Objective}

The primary goal of Snapshot 2 is to implement an interactive map and
location search system that allows users to monitor AQI data across
multiple regions.

\subsection{Feature Tasks}

\begin{enumerate}
    \item Integrate Mapbox GL JS
    \item Configure react-native-maps
    \item Build AQI heat-map overlays
    \item Implement Google Places autocomplete
    \item Add save location functionality
    \item Add reverse geocoding support
    \item Update mobile navigation system
\end{enumerate}

\subsection{New Dependencies}

\begin{longtable}{|p{5cm}|p{3cm}|p{5cm}|}
\hline
\textbf{Tool} & \textbf{Version} & \textbf{Purpose} \\
\hline
Mapbox GL JS & 3.1 & Interactive web maps \\
\hline
react-native-maps & 1.10 & Mobile map rendering \\
\hline
Google Places API & v1 & Location autocomplete \\
\hline
Turf.js & 6.5 & Geospatial calculations \\
\hline
React Navigation & 6.1 & Mobile navigation system \\
\hline
\end{longtable}


\subsection{TestRail \& Objectives}
TestRail Suite Added: Map \& Location Module (Suite ID: S-102)
\begin{itemize}
    \item TC-201: Verify Mapbox rendering on web
    \item TC-202: Validate AQI heat-map overlay accuracy
    \item TC-203: Test GPS location detection and reverse geocoding
    \item TC-204: Verify saved locations persistence after app restart
\end{itemize}

\subsection {Updated Workflow}
Sprint 3: Map integration and tile layer + Sprint 4: Location search, saved locations, QA pass
Added code-review checklist item: verify map API key is not committed to repository


\subsection{Documentation Updates}

The following documents will be revised:

\begin{itemize}
    \item Software Design Document
    \item Software Requirements Specification
    \item User Manual
    \item Workflow Diagram
\end{itemize}

\subsection{Expected Deliverables}

\begin{itemize}
    \item Interactive map system
    \item Multi-location support
    \item Search and geolocation features
    \item Updated project documentation
\end{itemize}
\newpage

\section{Snapshot 3 - Checkpoint 2 Reflection}

Snapshot 2 introduced interactive maps, location search,
and persistent saved locations functionality.

\subsection{New Feature Objective}

Snapshot 3 focuses on implementing push notifications and
a health advisory engine for AQI threshold monitoring.

\subsection{Feature Tasks}

\begin{enumerate}
    \item Integrate Firebase Cloud Messaging
    \item Configure APNs notifications
    \item Create notification preferences screen
    \item Build AQI scheduler service
    \item Implement Health Advisory Engine
    \item Build alerts feed interface
    \item Add notification deep-linking
\end{enumerate}

\subsection{New Dependencies}

\begin{longtable}{|p{5cm}|p{3cm}|p{5cm}|}
\hline
\textbf{Library/Tool} & \textbf{Version} & \textbf{Purpose} \\
\hline
Expo Notifications & 0.27 & Push notifications \\
\hline
Firebase Admin SDK & 12.0 & Notification dispatching \\
\hline
node-cron & 3.0 & Scheduled jobs \\
\hline
Bull Queue & 4.12 & Notification queue system \\
\hline
Redis & 7.2 & Queue backend \\
\hline
Joi & 17.12 & API validation \\
\hline
\end{longtable}

\subsection{TestRail \& Snapshot Objectives}
TestRail Suite Added: Notifications \& Advisory Engine (Suite ID: S-103)
\begin{itemize}
    \item TC-301: Verify Firebase Cloud Messaging delivery
    \item TC-302: Test AQI threshold logic and advisory generation
    \item TC-303: Validate quiet hours notification suppression
    \item TC-304: Verify advisory feed rendering and alert history display
\end{itemize}

\subsection {Updated Workflow}
Sprint 5: Notification infrastructure + Sprint 6: Health Advisory Engine, alerts UI, QA
Added Redis and Bull to docker-compose.yml for local dev environment

\subsection{Documentation Updates}

Updated documents include:

\begin{itemize}
    \item SDD notification architecture
    \item SRS notification requirements
    \item User manual alert configuration
    \item Workflow diagrams
\end{itemize}

\subsection{Expected Deliverables}

\begin{itemize}
    \item Fully operational notification system
    \item Health advisory recommendation engine
    \item Alerts feed interface
    \item Updated technical documentation
\end{itemize}
\newpage
\section{Snapshot 4 - Final Reflection}

AirWatch successfully delivered a cross-platform weather
and air quality monitoring platform featuring:

\begin{itemize}
    \item Real-time AQI monitoring
    \item Interactive heat-map visualization
    \item Push notifications
    \item Health advisory recommendations
\end{itemize}

\subsection{Final Development Goals}

\begin{enumerate}
    \item Optimize performance
    \item Implement dark mode
    \item Conduct accessibility audits
    \item Finalize deployment pipelines
    \item Complete regression testing
    \item Prepare production release
\end{enumerate}

\subsection{New Dependencies}

\begin{longtable}{|p{5cm}|p{3cm}|p{5cm}|}
\hline
\textbf{Tool} & \textbf{Version} & \textbf{Purpose} \\
\hline
react-native-dark-mode & 1.1 & Dark mode support \\
\hline
Lighthouse CI & 0.13 & Accessibility auditing \\
\hline
Fastlane & 2.219 & Deployment automation \\
\hline
Sentry & 7.99 & Crash monitoring \\
\hline
\end{longtable}

\subsection{Future Work}

Future planned enhancements include:

\begin{itemize}
    \item Wearable device integration
    \item AI-powered AQI forecasting
    \item Community pollution reporting
    \item Offline mode support
    \item Multi-language support
    \item Home screen widgets
\end{itemize}

\subsection{TestRail Reflection}
Total Test Cases Written: 92  |  Executed: 92  |  Passed: 86  |  Failed: 4  |  Blocked: 2

\begin{itemize}
    \item Most critical failures involved dark mode UI inconsistencies and AQI heat-map lag on older devices
    \item Blocked issues included incomplete APNs production certificate validation during final iOS testing
    \item Accessibility testing identified missing screen reader labels and insufficient contrast ratios
    \item Overall pass rate: 93\% --- meeting the project release threshold
\end{itemize}

\subsection{Release Workflow}

\begin{enumerate}
    \item Final QA verification
    \item Accessibility review
    \item Production build generation
    \item TestFlight and Play Store testing
    \item Version tagging (v1.0.0)
    \item Public release deployment
\end{enumerate}

\end{document}